
\textbf{Verification-in-Loop Systems.} AutoSpec+ demonstrated iterative refinement with LLMs and Frama-C achieves 96\% proof ratio on 604 C programs, establishing verification-in-loop as state-of-practice. LeanDojo showed theorem proving benefits from proof-assistant feedback in Lean. Our work extends these systems by decomposing why iteration works---quantifying information gradient across feedback dimensions ($\beta =12.49, R^{2}=0.89)$---and how to generalize via cross-verifier semantic normalization.

AutoSpec+ uses natural language error messages and proof-aware decomposition but does not decompose feedback structure into reusable dimensions or quantify information value. When a precondition fails, AutoSpec+ returns unstructured messages. Our approach extracts three dimensions: (1) Witness: concrete counterexample, (2) Structure: semantic category, (3) Dependency: causal chain. This structured extraction enables cross-verifier transfer and quantitative analysis.

\textbf{LLM-Guided Formal Methods.} PropertyGPT uses retrieval-augmented generation to achieve 80\% recall for smart contract property generation from natural language. Our approach provides complementary structured signal: PropertyGPT uses external knowledge (retrieved examples), we use internal constraints (verifier feedback). These are additive---RAG provides domain patterns, feedback provides program-specific constraints.

\textbf{Error Taxonomies and Cross-Tool Translation.} FormalRx introduced 28-category error taxonomy for proof assistants (Lean, Coq), demonstrating error classification generalizes across theorem provers. Building on FormalRx's cross-tool taxonomy approach, we demonstrate that SMT-based program verifiers enable a minimal 8-primitive taxonomy (vs 28 for proof assistants), achieving 100\% coverage across Frama-C, Dafny, Why3. This minimalism reflects the narrower semantic core of SMT-based verifiers compared to proof assistants.

\textbf{Mutation Testing for Specification Quality.} Mutation testing traditionally validates test suite quality. We apply mutation testing to LLM-synthesized formal specifications to address concerns about ''spec washing'' (trivial or vacuous specifications). Our results (63.3\% kill rate matching 60\% gold baseline) provide evidence that synthesized specs are semantically meaningful.
